\documentclass{standalone}
%% Show/hide steps for standalone work — edit this line:
\makeatletter\AtBeginDocument{\@buildstepset{1,2}}\makeatother
\usepackage{fontawesome5}
\usepackage{hsbcommon}
\begin{document}
\begin{tikzpicture}[remember picture]

        \tikzset{every node/.style={font=\sffamily}}

    \setcounter{pcell}{0}
    \matrix (M) [
        matrix of nodes,
        column sep=-\pgflinewidth,
        row sep=-\pgflinewidth,
        nodes={anchor=center, draw=none, inner sep=0pt},
        column 1/.style={nodes={minimum width=5em, align=right,color=red}},
        row 1/.style={nodes={minimum height=2em, align=center,color=blue}},
    ]{
             & \strut swerve              & \strut straight            \\
    swerve     & \payoff{\bdice{2}}{\rdice{2}}   & \payoff{\bdice{6}}{\rdice{1}}   \\
    straight     & \payoff{\bdice{1}}{\rdice{6}}     & \payoff{\bdice{}}{\rdice{}}   \\
    };


    % Alice: centered between column 2 and 3 header cells, above row 1
    \node[anchor=south,blue] at ($(M-1-2)!0.5!(M-1-3)$) [yshift=1em] {Alice};
    
    % Bob: rotated, centered between row 2 and 3, left of column 1  
    \node[anchor=south, rotate=90,red] at ($(M-2-1)!0.5!(M-3-1)$) [yshift=3em] {Bob};
    


    
    % % Now address any number directly:
    % \draw[thick] (AA-al) circle (10pt);   % circle Alice's 10
    % \draw[->]    (AB-al) -- (AA-al);      % arrow from Alice's 8 to 10


    \ifstep{1}{
        \draw[->,>=latex,thick,blue]  (AA-al) -- (AB-al) ;      
        % \draw[->,>=latex]    (BB-al) -- (AB-al);     
        % \draw[->,>=latex]    (BA-al) -- (AA-al);     
        \draw[<-,>=latex,thick,blue]    (BA-al) -- (BB-al);     
  
    }
    \ifstep{2}{
        \draw[<-,>=latex,red,thick]    (AB-bo) -- (BB-bo);     
        \draw[<-,>=latex,red,thick]    (BA-bo) -- (AA-bo);     
  
    }
   

  \ifstep{3}{
    \begin{scope}[on background layer]
        \fill[green!10] (M-2-3.north west) rectangle (M-2-3.south east);
        \fill[green!10] (M-3-2.north west) rectangle (M-3-2.south east);
        \fill[yellow!10] (M-2-2.north west) rectangle (M-2-2.south east);
        \fill[red!10] (M-3-3.north west) rectangle (M-3-3.south east);
    \end{scope}
    }






\end{tikzpicture}
\end{document}